\begin{abstract}
\noindent PhishVN-SMS is a \emph{pilot} corpus of Vietnamese SMS messages labelled in three
classes (\texttt{legitimate}, \texttt{spam}, \texttt{phishing}), contributed by consenting
volunteers from their own handsets. Only machine-sent messages were collected: scam and spam
broadcasts, and legitimate traffic from brandnames and service shortcodes; no person-to-person
message was taken. Personal data was redacted on the contributor's device before any text
reached the researchers, while links and broadcast amounts were kept as signal. Two annotators
labelled every message independently and disagreements were adjudicated; both original labels
ship unmodified with the agreement reported. Messages are grouped into templates by a stated,
reproducible computation, and the fixed 70/15/15 split is template-disjoint. A published
Vietnamese SMS corpus ships alongside, unchanged in label, as a held-out external benchmark. The
corpus holds \FILL{N} messages in \FILL{T} templates; the benchmark holds 2{,}676. It is a pilot
from one small group, not a sample of Vietnamese smishing. Released under \FILL{corpus licence}.
\end{abstract}
% \FILL{trim to <= 150 words once the counts are in}
